
\documentclass{article}
\usepackage{colortbl}
\usepackage{makecell}
\usepackage{multirow}
\usepackage{supertabular}

\begin{document}

\newcounter{utterance}

\centering \large Interaction Transcript for game `wordle{-}crazy\_withclue', experiment `medium\_frequency\_words\_clue\_no\_critic', episode 1 with LLP: Large Language Problems\_\_qwen3.5{-}9b\_\_llp{-}final.
\vspace{24pt}

{ \footnotesize  \setcounter{utterance}{1}
\setlength{\tabcolsep}{0pt}
\begin{supertabular}{c@{$\;$}|p{.15\linewidth}@{}p{.15\linewidth}p{.15\linewidth}p{.15\linewidth}p{.15\linewidth}p{.15\linewidth}}
    \# & \multicolumn{2}{c}{Player} && \multicolumn{2}{c}{Game Master} \\
    \hline

    \theutterance \stepcounter{utterance}  
    & & & \multicolumn{4}{p{0.6\linewidth}}{
        \cellcolor[rgb]{0.9,0.9,0.9}{
            \makecell[{{p{\linewidth}}}]{
                \texttt{\tiny{[P1$\langle$GM]}}
                \texttt{You are a language wizard who likes to guess words by using the given rules.} \\
\\ 
\texttt{Welcome to Wordle! You have six attempts to guess the target word, a valid English word of five lowercase letters (a{-}z). Please use the tags "explanation:" and "guess:" to provide a concise explanation for each guess.} \\
\\ 
\texttt{To help you make an informed guess, you will receive a clue for the word, such as} \\
\texttt{clue: snowy white.} \\
\\ 
\texttt{Here is an example guess based on the clue:} \\
\texttt{explanation: In the fairy tail Snow White, the girl is killed because she eats a poisoned apple. And the word apple has 5 letters.} \\
\texttt{guess: apple} \\
\\ 
\texttt{After each guess, your answer will be validated, and you will receive feedback indicating which letters are correct (yellow), which letters are correct but in the wrong position (black), and which letters are incorrect (purple). This feedback can be useful in determining which letters to include or exclude in your next guess.} \\
\\ 
\texttt{For example, the feedback for "apple" might be:} \\
\texttt{guess\_feedback: a<black> p<black> p<yellow> l<black> e<purple>} \\
\\ 
\texttt{The explanation should contain details about how both the guess\_feedback and clue work together in tandem to formulate a new guess.} \\
\\ 
\texttt{Do not generate any other text.} \\
\\ 
\texttt{Let's begin with your first guess.} \\
\\ 
\texttt{clue: logbook} \\
            }
        }
    }
     \\ \\

    \theutterance \stepcounter{utterance}  
    & \multicolumn{4}{p{0.6\linewidth}}{
        \cellcolor[rgb]{0.561,0.737,0.561}{
            \makecell[{{p{\linewidth}}}]{
                \texttt{\tiny{[P1$\rangle$GM]}}
                \texttt{explanation: A logbook is a record of events, often used in ships or aircraft. The word "log" is directly related to a logbook, and "log" can be extended to "logbook" itself. A 5{-}letter word related to logging or records is "logos" (meaning word or reason) or "logos" could be a stretch. More directly, "logbook" contains "book" and "log". A common 5{-}letter word associated with records or logs is "diary" (a personal logbook). Another possibility is "ledger" (a record book), but that's 6 letters. "Logos" is 5 letters but less direct. "Logos" means word or reason, not specifically a logbook. "Logos" might be too abstract. Let me think of words that mean a record or log: "diary" is 5 letters and is a type of logbook. "Logos" is 5 letters but not quite right. "Logos" could be interpreted as a collection of words or records. Actually, "logos" is more about meaning or word. A better guess might be "diary" as it's a personal logbook. Or "logos" as a word related to logs. Let me go with "diary" as it's a direct synonym for a personal logbook.} \\
\\ 
\texttt{guess: diary} \\
            }
        }
    }
    & & \\ \\

    \theutterance \stepcounter{utterance}  
    & & & \multicolumn{2}{p{0.3\linewidth}}{
        \cellcolor[rgb]{0.9,0.9,0.9}{
            \makecell[{{p{\linewidth}}}]{
                \texttt{\tiny{[GM$|$GM]}}
                \texttt{attempts = 1} \\
\texttt{target\_word = diary} \\
\texttt{target\_word\_clue = logbook} \\
\texttt{guess = diary} \\
\texttt{guess\_feedback = d<yellow> i<yellow> a<yellow> r<yellow> y<yellow>} \\
            }
        }
    }
    & & \\ \\

    \theutterance \stepcounter{utterance}  
    & & & \multicolumn{2}{p{0.3\linewidth}}{
        \cellcolor[rgb]{0.9,0.9,0.9}{
            \makecell[{{p{\linewidth}}}]{
                \texttt{\tiny{[GM$|$GM]}}
                \texttt{game\_result = WIN} \\
            }
        }
    }
    & & \\ \\

\end{supertabular}
}

\end{document}
